\chapter{Programacion entera}

\section{Motivacion}
La programacion entera extiende la programacion lineal al exigir que algunas o
todas las variables adopten valores enteros. Esta restriccion permite modelar
decisiones indivisibles como abrir instalaciones, seleccionar proyectos,
asignar turnos o producir lotes discretos \cite{nemhauser1988integer,wolsey1998integer}.

\section{Modelo general}
Un problema mixto entero lineal puede escribirse como
\[
  \min \; c^\top x + d^\top y
  \quad \text{sujeto a} \quad
  Ax + By \geq b,\;
  x \in \R_+^n,\;
  y \in \Z_+^p.
\]
La dificultad computacional crece porque la region factible deja de ser
convexa en el sentido continuo: la relajacion lineal puede proponer soluciones
fraccionarias que no son implementables.

\section{Ejemplo QSB: requerimientos minimos discretos}
El ejemplo \texttt{manual/examples/ilp\_sample.json} minimiza el costo de dos
actividades enteras:
\[
\begin{aligned}
  \min \quad & 2.5X_1 + 2X_2 \\
  \text{s.a.}\quad & 6X_1 + 3X_2 \geq 200,\\
  & 3X_1 + 5X_2 \geq 180,\\
  & X_1, X_2 \in \Z_+.
\end{aligned}
\]

Comando reproducible:
\begin{lstlisting}
swift run qsb solve-json-ilp manual/examples/ilp_sample.json
\end{lstlisting}

\begin{table}[h]
\centering
\caption{Solucion del ejemplo de programacion entera}
\begin{tabular}{lr}
\toprule
Cantidad & Valor \\
\midrule
\(X_1\) & 22 \\
\(X_2\) & 23 \\
Valor objetivo & 101 \\
\bottomrule
\end{tabular}
\end{table}

\section{Notas de verificacion}
La solucion satisface las restricciones:
\[
6(22)+3(23)=201 \geq 200,\qquad 3(22)+5(23)=181 \geq 180.
\]
La cercania de ambas restricciones a sus requerimientos muestra el equilibrio
tipico entre costo minimo e integridad de las decisiones.
